% =============================================================================
% SECTION IV — Cross-Strategy Interdependencies
% =============================================================================


% ── SEMAPHORE ────────────────────────────────────────────────────────────────

The preceding sections have established that all three strategies
generate alpha that survives both classical benchmark and data-driven
factor controls, and that this alpha rises during periods of
intermediary funding stress. A natural question follows: are the
three mispricings linked to one another, and if so, what economic
mechanism ties them together? The slow-moving capital literature
introduced above---\citet{duffie2010presidential},
\citet{brunnermeier2009market}, and \citet{mitchell2012arbitrage}---yields
five testable predictions about how arbitrage mispricings behave when they
draw on a common pool of slow-moving intermediary capital:
(P1)~mispricings serviced by a common pool of intermediary capital
should co-move \citep{duffie2010presidential}; (P2)~the co-movement
should intensify when that capital is scarce
\citep{duffie2010presidential, brunnermeier2009market};
(P3)~mispricing changes in one strategy should contain information for
predicting those of others, with effects propagating asymmetrically
along the liquidity gradient \citep{mitchell2012arbitrage, fleckenstein2014tips};
(P4)~the persistence of mispricing levels should reflect the
funding-dependence of each arbitrage trade, with strategies that
require continuous balance-sheet funding displaying the longest
half-lives \citep{boyarchenko2016trends}; and (P5)~the mispricings
should respond to the shadow cost of intermediary funding
\citep{duffie2010presidential, brunnermeier2009market}.

A feature of our three-strategy setting gives these tests their
identifying power. All three arbitrages are run by the same type of
leveraged institutional arbitrageur, but the instruments differ in who
holds them: the CDS--Bond and iTraxx legs trade among institutions,
whereas the BTP~Italia inflation-linked bond is held predominantly by
retail investors. This heterogeneity in the holder base is the source
of identification we exploit throughout---it lets us separate the
slow-moving-capital mechanism, under which co-movement depends on
sharing a capital pool, from a common macroeconomic driver, under which
all three mispricings would move together regardless of who holds the
instruments.

We devote one subsection to each prediction, stating in each case the
prediction, the test, and the verdict where the evidence is presented.
Section~\ref{sec:rq3_construction} first defines the mispricing series
and the stress regimes. Sections~\ref{sec:rq3_p1}--\ref{sec:rq3_p5}
then test P1 through P5 in turn. Section~\ref{sec:rq3_discussion}
interprets the evidence as a whole, explains why these arbitrage
opportunities persist in post-crisis European credit markets, and
collects the verdicts in a summary mapping.


% ── 6.1 MISPRICING CONSTRUCTION ──────────────────────────────────────────────

\subsection{Mispricing Construction and Stress Regimes}
\label{sec:rq3_construction}

Following \citet{duffie2010presidential}, we work with mispricing
magnitudes $m_{i,t}$ rather than strategy returns, because it is the
level of the mispricing---the distance of the basis from its
no-arbitrage value---that measures the degree to which capital is
failing to close the gap. For each strategy, $m_{i,t}$ is defined as
the absolute value of the basis, sign-adjusted so that an increase
$\Delta m > 0$ denotes a widening of the mispricing for all three
strategies. The BTP~Italia mispricing is the cross-sectional mean of the absolute
basis across all outstanding issues with at least six months to
maturity; the CDS--Bond mispricing is the cross-sectional mean of the absolute 
value of all negative single-name bases in the iTraxx Main universe with at least
six months to maturity on each day; and the iTraxx Skew mispricing is the
cross-sectional mean of the absolute skew across the four on-the-run
sub-indices. All three series are aggregated to monthly frequency.
Market conditions are classified into three regimes---LOW, MEDIUM, and
HIGH---by the iTraxx Europe Main 5-year spread, at the $60$ and
$100$~basis-point cutoffs used throughout the paper for the
transaction-cost tiers (Appendix~\ref{app:fees}) and the stress-regime
analyses of Sections~\ref{sec:benchmarks}
and~\ref{sec:factor_selection}.


% ── 6.2 P1: CROSS-MARKET CO-MOVEMENT ─────────────────────────────────────────

\subsection{P1: Cross-Market Co-Movement}
\label{sec:rq3_p1}

The first prediction is that mispricings serviced by a common pool of
intermediary capital should co-move: when that capital is withdrawn,
every arbitrage drawing on it loses the balance-sheet capacity that
compresses its mispricing, so they widen together. We test this with
pairwise Pearson correlations of the monthly $\Delta m$ series
(Table~\ref{tab:rq3_unified_corr}).

The prediction holds for the pair whose underlying instruments are all
institutionally held and traded---the CDS--Bond basis and the iTraxx
skew, hereafter the ``institutional pair.'' Their mispricing changes
correlate at $+0.19$ unconditionally, significant at the $10\%$ level,
the positive sign expected when both draw on the same pool of
intermediary capital.

The pairs involving BTP~Italia behave differently. The CDS--Bond
versus BTP~Italia correlation is $-0.30$, significant at the $1\%$
level, and the BTP~Italia versus iTraxx correlation is $-0.10$,
indistinguishable from zero.

This does not break P1, which predicts co-movement only for trades
that share the pool---and the retail-held BTP~Italia does not. The negative
sign is itself informative: it is the first empirical signal of the
holder-base segmentation we exploit for identification, since a
common macroeconomic driver would move all three mispricings together
regardless of who holds them. P1 is confirmed for the institutional
pair; we take up the economics of the negative sign in
Section~\ref{sec:rq3_discussion}.

% ── 6.3 P2: STRESS AMPLIFICATION ─────────────────────────────────────────────

\subsection{P2: Stress Amplification}
\label{sec:rq3_p2}

The second prediction sharpens the first: the co-movement should not
merely be present but should intensify as intermediary capital becomes
scarce. A first reading comes from the regime-conditional correlations
(Table~\ref{tab:rq3_unified_corr}, plotted by regime in
Figure~\ref{fig:rq3_regime_bar}): within the institutional pair the
correlation rises monotonically from $-0.21$ in LOW to $+0.32$ in
HIGH, with the LOW-to-HIGH movement significant at
the $1\%$ level under a Fisher $z$-test ($p=0.005$), and the \citet{forbes2002no}
correction for the heteroskedastic increase in mispricing volatility
under stress leaves it significant at the $10\%$ level ($p=0.070$)---so the
intensification is a genuine increase in the cross-market link, not a
volatility artefact. But the regime split treats stress as a discrete
state and cannot say whether transmission rises continuously with it,
nor separate the incremental effect of stress from the average level
of transmission. We therefore test the prediction
parametrically with an interaction regression:
\begin{equation}\label{eq:interaction}
  \Delta m_{i,t} \;=\; \alpha \;+\; \beta_0\, \Delta m_{j,t}
  \;+\; \beta_1\, (\Delta m_{j,t} \times z_t)
  \;+\; \gamma\, z_t \;+\; \varepsilon_t,
\end{equation}
where $z_t$ is the contemporaneous standardised iTraxx Main 5-year
spread, a proxy for intermediary funding stress; the state enters
contemporaneously because the object here is amplification, not
predetermined conditional performance. The coefficient $\beta_1$ captures the
incremental transmission during stress: $\beta_1 > 0$ means that a
given shock to $\Delta m_j$ has a larger effect on $\Delta m_i$ when
balance sheets are strained. The level term $\gamma\,z_t$ absorbs the
direct widening of $\Delta m_i$ under stress, so that $\beta_1$
isolates amplification rather than the widening common to both series. The theory is explicit on the mechanism:
\citet{duffie2010presidential} notes that a depleted intermediary is
``less able or predisposed to absorb supply and demand shocks,'' and
\citet{brunnermeier2009market} show that the destabilising effects of
funding constraints are ``largest when traders are near their
constraint.''

Table~\ref{tab:rq3_interaction} reports the results. Within the
institutional pair the interaction term is correctly signed and, in
the $\textrm{CDS--Bond}\rightarrow\textrm{iTraxx}$ direction,
significant at the $5\%$ level ($\beta_1 = +0.08$). Regime split and
interaction point the same way, and P2 is confirmed for the
institutional pair on both tests.

For the retail-held BTP~Italia the sign flips. The CDS--Bond versus
BTP~Italia correlation moves from $-0.13$ in LOW to $-0.47$ in HIGH,
significantly negative in the MEDIUM and HIGH states at the $1\%$ and
$5\%$ levels; the BTP~Italia versus iTraxx correlation turns from
$+0.15$ to $-0.14$, significant in none. Neither pair shows a
significant LOW-to-HIGH change, so the interaction regression, which
does not split the sample, carries the result
($\beta_{1} = -0.64$ for
$\textrm{CDS--Bond}\rightarrow\textrm{BTP~Italia}$, significant at the
$1\%$ level): stress does not amplify a positive link here but deepens
a negative one. The bivariate GARCH--DCC estimates
(Appendix~\ref{app:dcc}, Figure~\ref{fig:app_dcc}) show this is no
artefact of regime averaging: both correlations are negative in every
month, averaging $-0.29$ against the CDS--Bond basis and $-0.10$
against iTraxx. This is the mirror image of the institutional
pattern---exactly what the holder-base segmentation of
Section~\ref{sec:rq3_discussion} predicts.

% ── 6.4 P3: LEAD-LAG TRANSMISSION ────────────────────────────────────────────

\subsection{P3: Lead--Lag Transmission}
\label{sec:rq3_p3}

The third prediction concerns timing: if capital moves slowly across
markets, a mispricing change in one strategy should help predict
future changes in another, and the transmission should run from more
liquid to less liquid markets as capital is reallocated across the
dealer network. We test predictability with spanning regressions in
the design of \citet{fleckenstein2014tips},
\begin{equation}\label{eq:spanning}
  \Delta m_{i,t} \;=\; \alpha
  \;+\; \sum_{j \neq i} \sum_{l=0}^{L} \beta_{j,l}\, \Delta m_{j,t-l}
  \;+\; \varepsilon_t,
\end{equation}
with $L = 2$ lags and Newey--West HAC standard errors. The objects of
interest are the lagged coefficients: a significant loading on
$\Delta m_{j,t-1}$ means that last month's mispricing in strategy $j$
predicts this month's in strategy $i$---a lead--lag relation. The
contemporaneous terms enter as controls, so that these lagged loadings
do not absorb the same-month co-movement of Section~\ref{sec:rq3_p1}.

Table~\ref{tab:rq3_spanning} reports the results. The lead--lag
evidence rests on a single loading: one-month-lagged iTraxx changes
predict the widening of the CDS--Bond mispricing, at $+0.71$ and the
$10\%$ level. The contemporaneous controls are significant in both
directions between the CDS--Bond basis and BTP~Italia, as the
co-movement of Section~\ref{sec:rq3_p1} implies. A small loading of
twice-lagged BTP~Italia
changes in the iTraxx equation ($+0.03$, $10\%$) runs the other way;
the causality tests below do not corroborate it.

The spanning regressions omit each series' own lags, so a predictive
loading cannot be told apart from persistence in the variable being
predicted. We therefore embed the three $\Delta m$ series in a
first-order VAR, which conditions on that history, and test for Granger
causality (Table~\ref{tab:rq3_granger_multilag}). The liquidity pecking
order is sharp. The iTraxx index Granger-causes both the CDS--Bond basis
($p=0.011$) and the BTP~Italia basis ($p=0.007$), and both links
survive at two and three lags. The CDS--Bond basis in turn
Granger-causes BTP~Italia at the $10\%$ level, completing the chain.
Neither direction reverses: BTP~Italia does not Granger-cause either
institutional mispricing at conventional levels, nor does CDS--Bond
Granger-cause iTraxx. Shocks thus originate in the most liquid,
standardised index market, where dealer inventory adjusts fastest, and reach both
the less liquid cash basis and the retail-held inflation-linked basis
within a month. The generalised impulse responses of \citet{pesaran1998generalized},
invariant to the ordering of variables in the VAR, are consistent
with the same picture (Figure~\ref{fig:rq3_girf}). A
one-standard-deviation shock to the iTraxx index produces a positive
response in the CDS--Bond basis and a negative response in the
BTP~Italia basis, both largest on impact and significant there, and
both decaying to zero within two to three months. Shocks in the reverse
direction---originating in either of the less liquid bases and propagating back
to the iTraxx skew---are far smaller and indistinguishable from zero
beyond impact. The asymmetric pattern matches the forced-deleveraging
mechanism of \citet{mitchell2012arbitrage}, by which intermediaries
adjust their most liquid positions first when funding contracts and
their less liquid positions only with a delay.
P3 is confirmed: predictability, direction, and dynamic response all run
from the liquid index to the less liquid bases.

% ── 6.5 P4: PERSISTENCE AND FUNDING-DEPENDENCE ───────────────────────────────

\subsection{P4: Persistence and Funding-Dependence}
\label{sec:rq3_p4}

The fourth prediction links the persistence of each mispricing to how
funding-dependent its trade is: the more a strategy relies on
continuously financing a cash position on a dealer balance sheet, the
longer a shock to its mispricing should take to decay, because the
capital that closes it arrives only slowly. \citet{duffie2010presidential}
makes this quantitative---the half-life of the reversal that follows a
shock is decreasing in the mean rate at which arbitrage capital reaches
the trade---so half-lives should track funding frictions across
strategies. We estimate first-order
autocorrelations of the mispricing levels,
\begin{equation}\label{eq:ar1}
  m_{i,t} = c_i + \phi_i\,m_{i,t-1} + u_{i,t},
\end{equation}
with Newey--West standard errors, and read off the
implied half-life $h_i^{*} = \ln(0.5)/\ln(\phi_i)$.

Table~\ref{tab:rq3_ar1_persistence} reports the estimates. The
CDS--Bond basis is the most persistent ($\phi = 0.79$, half-life
$\approx 3.0$ months), with the BTP~Italia and the iTraxx skew
essentially identical at $\phi \approx 0.73$ and half-life
$\approx 2.2$ months, all three coefficients significant at the
$1\%$ level. The CDS--Bond basis is the only one of the three that is
not cash-neutral: closing it means buying the bond and financing it for
as long as the position is held; the CDS leg hedges the credit risk but
requires no cash of its own. When funding is scarce the trade cannot be
put on, and the basis stays wide until fresh capital arrives---roughly
one-third longer than for the other two, neither of which requires
that: the BTP~Italia basis is long one cash bond against another, so
the short leg funds the long, and the iTraxx skew is a pure derivative
trade with no bond to finance. On funding grounds the two are
equivalent, though they sit at opposite ends of the liquidity ordering
of Section~\ref{sec:rq3_p3}.

The diagonal panels of Figure~\ref{fig:rq3_girf} show the same process
on the other timescale: an innovation to $\Delta m$ is absorbed within
the month, while the level it shifts unwinds over the horizon the
AR(1) measures. A sharp reaction and a slower recovery is what
\citet{duffie2010presidential} associates with slowly arriving capital.

P4 is confirmed on the one distinction the three strategies allow. The
trade that must be funded is a third more persistent than the two that
need not be, and those two are indistinguishable in persistence as the
funding criterion implies.


% ── 6.6 P5: FUNDING CONDITIONS AS THE DRIVER ───────────────────────────────────────

\subsection{P5: Funding Conditions as the Driver}
\label{sec:rq3_p5}

The fifth prediction concerns what moves the mispricings.
\citet{duffie2010presidential} places the availability of intermediary
capital among the variables that price assets---``equilibrium
asset-pricing fundamentals include state variables determining the
immediate and future availability of capital''---and
\citet{brunnermeier2009market} make the channel precise: a security's
illiquidity is ``the product of its margin and the shadow cost of
funding.'' If slow-moving capital is what keeps these bases open,
measures of that shadow cost should move them. We regress each
$\Delta m$ series on four proxies
(Table~\ref{tab:rq3_pc1_funding})---the \citet{he2017intermediary}
ratio of primary-dealer equity to assets, the Libor--OIS spread, the
change in the five-year CDS spread of a basket of five European prime
brokers, and the \citet{amihud2015pricing} price-impact measure---on
the months for which all three mispricings are available, so that the
coefficients can be compared across strategies.

Funding conditions move all three mispricings, and not in the same
direction. The Libor--OIS spread is significant in every equation and
reverses sign across the holder bases: $+16.63$ on the CDS--Bond basis
and $+3.85$ on the iTraxx skew, both at the $1\%$ level, against
$-15.99$ on BTP~Italia at the $5\%$ level. Illiquidity repeats the
pattern where it enters, $+6.93$ on the CDS--Bond basis at the $5\%$
level against $-20.51$ on BTP~Italia at the $1\%$ level. Dealer credit
risk loads on the CDS--Bond basis alone ($+0.11$, at the $10\%$
level)---the one trade of the three that has to be financed outright,
and the one whose persistence Section~\ref{sec:rq3_p4} traced to that
requirement. The \citet{he2017intermediary} ratio is insignificant
throughout; a quarterly capital ratio interpolated to monthly frequency
is a coarse instrument for a constraint that binds week to week.

The sign reversal is the institutional-versus-retail pattern of
Sections~\ref{sec:rq3_p1} and~\ref{sec:rq3_p2}, now appearing in the
response to the driver itself, and it is what identifies the channel:
an aggregate risk premium would move two claims on the same cash flows
in the same direction whoever holds them, whereas a margin constraint
binds only where the holder is levered. P5 is confirmed---the shadow
cost of intermediary funding drives the mispricings---and we take up
the economics of the negative sign in
Section~\ref{sec:rq3_discussion}.


% ── 6.7 DISCUSSION ───────────────────────────────────────────────────────────

\subsection{Discussion: Why These Arbitrages Persist}
\label{sec:rq3_discussion}

\paragraph{Holder-Base Heterogeneity as the Source of Identification.}

The evidence of the preceding subsections is unified by the single
structural feature introduced at the outset: although all three
strategies are run by the same type of institutional arbitrageur, the
instruments differ in who holds them, and it is this heterogeneity that
identifies the mechanism. The CDS--Bond basis and the iTraxx skew trade
in instruments whose holder base is overwhelmingly
institutional---hedge funds, dealer proprietary desks, and leveraged
accounts---so that a spike in volatility inflates the Value-at-Risk of
the combined position, exhausts the desk's risk budget, and forces
simultaneous liquidation across both books \citep{brunnermeier2009market},
pushing both bases wider at once. The BTP~Italia is different: only its
inflation-linked leg is structurally retail-dominated. Introduced in
2012 with a primary auction reserved to households and a loyalty bonus
(\textit{premio fedelt\`{a}}) for buy-and-hold investors, and serving
as the simplest inflation hedge available to an Italian household, the
bond is held by preferred-habitat investors in the sense of
\citet{vayanos2021preferred}: they use no leverage, face no margin
calls, and do not unwind when dealer capital is scarce. The arbitrageur
of the BTP~Italia basis is itself institutional and faces the same
funding environment as the others---so the strategy is fully part of
the experiment, not a bystander to it---but the absence of forced
selling on the bondholder side keeps the basis out of the institutional
fire sales, so that stress drives it opposite to the institutional pair
rather than alongside it. This is what gives the design its identifying
power: a common macroeconomic factor would move all three mispricings
in the same direction regardless of who holds the instruments, whereas
the sign of the co-movement turns entirely on whether the holder base
is shared---the pattern documented under P1 and P2 and, in the response
to funding conditions themselves, under P5.

\paragraph{Why the Mispricings Persist.}

Why have these opportunities not been competed away? The answer lies in
structural limits to arbitrage that differ across strategies but share
a common root in the scarcity of intermediary capital. For the
BTP~Italia, the binding constraint is the segmentation itself:
institutional arbitrageurs can exploit the basis, but the retail
holders who form the natural counterparty neither monitor nor arbitrage
it \citep{vayanos2021preferred}. For the CDS--Bond basis, the limits
are those of \citet{bai2019cds} and \citet{mitchell2012arbitrage}: the
trade must fund the corporate bond on a dealer balance sheet, the repo
market for European investment-grade corporates is thin and subject to
recall risk, and the counterparty risk on the CDS leg consumes
regulatory capital---so capacity is procyclical and the mispricing
countercyclical, widening precisely when the trade is most attractive
but least feasible. For the iTraxx skew, the constraint is post-crisis
regulation: \citet{boyarchenko2016trends} show that the supplementary
leverage ratio and other Basel~III requirements raised the capital
charge on index-versus-constituent trades---as required leverage ratios
rose from roughly $2$--$3\%$ pre-crisis to $5$--$6\%$ post-crisis, the
return on equity fell by a factor of two to three---rendering them
``significantly less economical'' and producing ``potentially `new
normal' levels for both the CDS-bond and the CDX-CDS bases.'' The
persistence of alpha documented in
Sections~\ref{sec:benchmarks}~and~\ref{sec:factor_selection} is
therefore not a puzzle calling for a risk-based explanation; it is the
equilibrium of a market in which the cost of intermediary capital has
been permanently raised, and slow-moving capital, in the sense of
\citet{duffie2010presidential}, keeps mispricings open longer than
frictionless models would predict \citep{mitchell2007slow}.

\paragraph{Summary: The Five Predictions.}

\begin{enumerate}[label=\textbf{P\arabic*.}, leftmargin=*, itemsep=4pt]

\item Mispricings drawing on a common capital pool co-move. The
  institutional pair correlates at $+0.19$, while the retail-held
  BTP~Italia, outside that pool, correlates significantly negatively
  with the CDS--Bond basis at $-0.30$ unconditionally
  (Table~\ref{tab:rq3_unified_corr}). The sign turns on the holder
  base, exactly as the common-capital mechanism requires. Confirmed.

\item The co-movement intensifies in stress. Institutional co-movement
  rises monotonically from $-0.21$ in LOW to $+0.32$ in HIGH, a
  LOW-to-HIGH increase significant at the $1\%$ level under a
  Fisher $z$-test and at the $10\%$ level after the
  \citet{forbes2002no} adjustment, with the interaction term correctly
  signed and significant at the $5\%$ level
  (Tables~\ref{tab:rq3_unified_corr},~\ref{tab:rq3_interaction}).
  Confirmed.

\item Mispricing changes transmit along the liquidity gradient. The
  iTraxx index Granger-causes both other mispricings ($p=0.007$ and
  $0.011$), with impulse responses tracing the same liquid-to-illiquid
  path (Table~\ref{tab:rq3_granger_multilag},
  Figure~\ref{fig:rq3_girf}). Confirmed.

\item Persistence reflects funding-dependence. The CDS--Bond basis,
  the only strategy requiring continuous cash-bond financing on a
  dealer balance sheet, has a half-life about one-third longer than
  the iTraxx skew and the BTP~Italia
  (Table~\ref{tab:rq3_ar1_persistence}), which are equivalent on
  funding grounds and equivalent in persistence. Confirmed.

\item Funding conditions drive the mispricings, and split them by
  holder base. The Libor--OIS spread enters the CDS--Bond and iTraxx skew
  equations at $+16.63$ and $+3.85$ (both $1\%$) and BTP~Italia at
  $-15.99$ ($5\%$), with illiquidity repeating the pattern; dealer
  credit risk loads only on the one trade that must be funded
  (Table~\ref{tab:rq3_pc1_funding}). Confirmed.

\end{enumerate}

\noindent All five predictions find support in the data.